
%%%%%%%%%%%%%%%%%%%%%%%%%%%%%%%%%%%%%%%%
%           main body
%%%%%%%%%%%%%%%%%%%%%%%%%%%%%%%%%%%%%%%%

%-----------------------------------
%	TITLE PAGE
%-----------------------------------

\title[数学分析]{\LARGE \bfseries 第九章\quad 数项级数} % The short title appears at the bottom of every slide, the full title is only on the title page
\author[Cesàro可和级数] % Your institution as it will appear on the bottom of every slide, maybe shorthand to save space
{\Large \bfseries 级数的Cesàro求和法及Tauber型定理}
% \author{Singular Value Decomposition} % Your name
% \date{\today} % Date, can be changed to a custom date
\date{}

%------------------------------------------------

\begin{document}


%------------------------------------------------

\begin{frame}

\titlepage % Print the title page as the first slide

\end{frame}

%------------------------------------------------

% \begin{frame}
% \frametitle{内容提要} % Table of contents slide, comment this block out to remove it
% \tableofcontents % Throughout your presentation, if you choose to use \section{} and \subsection{} commands, these will automatically be printed on this slide as an overview of your presentation
% \end{frame}

%-----------------------------------
%	PRESENTATION SLIDES
%-----------------------------------

%------------------------------------------------

\begin{frame}{引言}

这一讲，我们来扩展“可求和”级数的范围，例如

\pause

$$\sum\limits_{n=1}^{\infty} (-1)^{n+1} = 1 - 1 + 1 - 1 + \cdots$$

\pause

这在通常意义下是一个发散的级数, 但我们很快会知道它是一个所谓的Cesàro 意义下可和的级数.

\vspace*{2em}
\pause

Cesàro 求和法, 实际上可以视作是某种“动态加权平均”的方法. 配合上Tauber型的定理, 可以在掌握数项级数某种意义上的前 $n$ 项均值变化趋势的前提下, 结合关于通项的某种“正则性”, 对级数本身的敛散性等性质做出判断.

\end{frame}

%------------------------------------------------

\section[定义]{Cesàro可和级数定义}

%------------------------------------------------

\begin{frame}{Cesàro意义下可和的级数}

\begin{Def}[Cesàro意义下可和的级数]
设 $s_n = \sum\limits_{k=1}^{n} a_k$ 是数项级数 $\sum\limits_{n=1}^{\infty} a_n$ 的前 $n$ 项和, 记 {\color{red}$\sigma_n = \dfrac{1}{n} \sum\limits_{k=1}^{n} s_k$}. 若 {\color{red}$\lim\limits_{n\to\infty} \sigma_n = A \in \mathbb{R}$}, 则称级数 $\sum\limits_{n=1}^{\infty} a_n$ 在Cesàro意义下可和, 或者更准确的说是在 $(c, 1)$ 的意义下可和, 并记为 $\sum\limits_{n=1}^{\infty} a_n = A~(c, 1).$
\end{Def}

\pause

$\sigma_n$ 实际上可以视作下面的加权平均:
$$
\begin{array}{r:cccccccccc}
  s_1 \phantom{++} & & a_1 \\
+ s_2 \phantom{++} & \phantom{+} + & a_1 & + & a_2 \\
+ s_3 \phantom{++} & \phantom{+} + & a_1 & + & a_2 & + & a_3 \\
\vdots \phantom{++} & & \vdots \\
+ s_n \phantom{++} & \phantom{+} + & a_1 & + & a_2 & + & a_3 & + & \cdots & + & a_n \\
\hdashline
\text{weight} \phantom{+} & & 1 & & \frac{n-1}{n} & & \frac{n-2}{n} & & \cdots & & \frac{1}{n}
\end{array}
$$

\end{frame}

%------------------------------------------------

\begin{frame}{Cesàro意义下可和级数的例子}

\begin{eg}
{\smaller[1]
考虑我们一开始提到的通常意义下发散的级数
$$\sum\limits_{n=1}^{\infty} a_n = \sum\limits_{n=1}^{\infty} (-1)^{n+1}  = 1 - 1 + 1 - 1 + \cdots,$$
容易算得
$$s_n = \sum\limits_{k=1}^{n} a_k = \begin{cases}
    1, & n = 2m - 1 \\
    0, & n = 2m
\end{cases} \quad m = 1, 2, \dots$$
进而有
$$\sigma_n = \sum\limits_{k=1}^n s_k = \begin{cases}
  \dfrac{m}{2m - 1}, & n = 2m - 1\\
  \dfrac{1}{2}, & n = 2m
  \end{cases} \quad m = 1, 2, \dots$$
}
$\sigma_n$ 关于 $n$ 有极限, 为 $\lim\limits \sigma_n = \dfrac{1}{2}.$ 于是这个级数在 $(c, 1)$ 意义下可和, 即有
$$\sum\limits_{n=1}^{\infty} a_n = \dfrac{1}{2}~(c, 1).$$
\end{eg}

\end{frame}

%------------------------------------------------

\section[性质]{Cesàro可和级数性质}

%------------------------------------------------

\begin{frame}{Cesàro意义下可和级数的性质}

一个自然的问题是：对于通常意义下收敛的级数 $\sum\limits_{n=1}^{\infty} a_n = A \in \mathbb{R},$ 它在 $(c, 1)$ 意义下是可和的吗？如果是, 那么它的和还是 $A$ 吗？

\pause
\vspace*{2em}

\begin{prop}
设级数 $\sum\limits_{n=1}^{\infty} a_n$ 在通常意义下收敛到有限实数 $A,$ 即 $\lim\limits_{n \to \infty} s_n = A,$ 那么级数 $\sum\limits_{n=1}^{\infty} a_n$ 是 $(c, 1)$ 可和的, 而且有 $\sum\limits_{n=1}^{\infty} a_n = A~(c, 1).$
\end{prop}

\end{frame}

%------------------------------------------------

\begin{frame}{Cesàro意义下可和级数的性质}

\begin{proof}
若 $\sum\limits_{n=1}^{\infty} a_n$ 在通常意义下收敛到有限实数 $A,$ 那么
$$A = \lim\limits_{n \to \infty} s_n = \lim\limits_{n \to \infty} \dfrac{n\sigma_n - (n-1)\sigma_{n-1}}{n - (n-1)}, $$
于是根据Stolz公式 (Stolz–Cesàro定理), 有
$$\lim\limits_{n \to \infty} \sigma_n = \lim\limits_{n \to \infty} \dfrac{n\sigma_n}{n} = \lim\limits_{n \to \infty} \dfrac{n\sigma_n - (n-1)\sigma_{n-1}}{n - (n-1)} = A.$$
这表明了 $\sum\limits_{n=1}^{\infty} a_n = A~(c, 1).$
\end{proof}

\end{frame}

%------------------------------------------------

\begin{frame}{$(c, r)$ 意义下可和的级数}

对于序列 $\{\sigma_k\},$ 可以考虑极限 $\lim\limits_{n\to\infty} \dfrac{1}{n} \sum\limits_{k=1}^n \sigma_k,$ 若此极限存在，则称级数 $\sum\limits_{n=1}^{\infty} a_n$ 是 $(c, 2)$ 意义下可和的级数:
$$
\begin{array}{r:cccccccccc}
  \sigma_1 \phantom{++} & & a_1 \\
+ \sigma_2 \phantom{++} & \phantom{+} + & a_1 & + & {\color{red}\frac{1}{2}}a_2 \\
\vdots \phantom{++} & & \vdots \\
+ \sigma_n \phantom{++} & \phantom{+} + & a_1 & + & {\color{red}\frac{n-1}{n}}a_2 & + & & \cdots & + & {\color{red}\frac{1}{n}}a_n \\
\hdashline
\text{weight} \phantom{+} & & 1 & & {\color{red}\frac{1}{n}\sum\limits_{k=2}^n \frac{k-1}{k}} & & & \cdots & & {\color{red}\frac{1}{n^2}}
\end{array}
$$

\pause
\vspace*{1em}

由此, 可以归纳地定义 $(c, r)$ 意义下可和的级数, $r \geqslant 1$ 为正整数.
$$\sigma_n^{(r)} = \dfrac{1}{n} \sum\limits_{k=1}^n \sigma_n^{(r-1)} \text{ 关于 $n$ 收敛}, ~~ (\text{约定 } \sigma_n^{(1)} = \sigma_n, \sigma_n^{(0)} = s_n).$$

\end{frame}

%------------------------------------------------

\begin{frame}{$(c, r)$ 意义下可和的级数}

将通常意义下收敛的级数称作是 $(c, 0)$ 意义下可和的级数, 那么根据之前我们证明的命题以及举的例子, 我们有如下的集合 (空间) 的真包含的关系:
% \begin{multline*}
%   \{ (c, 0) \text{可和级数} \} \subsetneqq \{ (c, 1) \text{可和级数} \} \subsetneqq \cdots \\ \subsetneqq \{ (c, r) \text{可和级数} \} \subsetneqq \{ (c, r+1) \text{可和级数} \} \subsetneqq \cdots
% \end{multline*}

$$
\begin{array}{c}
\{ (c, 0) \text{可和级数} \} \\
\rotatebox[origin=c]{-90}{$\subsetneqq$} \\
\{ (c, 1) \text{可和级数} \} \\
\rotatebox[origin=c]{-90}{$\subsetneqq$} \\
\vdots \\
\{ (c, r) \text{可和级数} \} \\
\rotatebox[origin=c]{-90}{$\subsetneqq$} \\
\{ (c, r+1) \text{可和级数} \} \\
\rotatebox[origin=c]{-90}{$\subsetneqq$} \\
\vdots
\end{array}
$$

\end{frame}

%------------------------------------------------

\section[Tauber型定理]{Tauber型定理}

%------------------------------------------------

\begin{frame}{Tauber型定理}

反过来, 我们有一些结论, 对 $(c, r+1)$ 可和的级数, 加上一些额外的限制, 能得到它落在子集, $\{ (c, r) \text{可和级数} \}$ 中. 这就是所谓的Tauber型的定理.

\pause

\vspace*{1em}

\begin{thm}[Tauber]
  若级数 $\sum\limits_{n=1}^{\infty} a_n = A~(c, 1),$ 且当 $n\to\infty$ 时 $a_n = \mathrm{o} \left( \dfrac{1}{n} \right),$ 那么该级数在通常意义下可和, 且 $\sum\limits_{n=1}^{\infty} a_n = A.$
\end{thm}

\end{frame}

%------------------------------------------------

\begin{frame}{Tauber型定理证明}

我们考察 $s_n - \sigma_n,$
\pause
有
\begin{align*}
  s_n - \sigma_n & = \sum\limits_{k=1}^n a_k - \sum\limits_{k=1}^n \left( 1 - \dfrac{k-1}{n} \right) a_k \\
  & = \dfrac{\sum\limits_{k=1}^n (k-1)a_k}{n}.
\end{align*}
\pause
由于 $n \to \infty$ 时有 $a_n = \mathrm{o}\left( \dfrac{1}{n} \right),$ 那么 $\lim\limits_{n \to \infty} n a_n = 0,$ 这也等价于
$$\lim\limits_{n \to \infty} (n - 1) a_n = 0.$$

\end{frame}

%------------------------------------------------

\begin{frame}{Tauber型定理证明}

于是对 $\dfrac{\sum\limits_{k=1}^n (k-1)a_k}{n}$ 使用Stolz公式有:
\begin{align*}
  \lim\limits_{n \to \infty} (s_n - \sigma_n) & = \lim\limits_{n \to \infty} \dfrac{\sum\limits_{k=1}^n (k-1)a_k}{n} \\
  & = \lim\limits_{n \to \infty} \dfrac{\sum\limits_{k=1}^n (k-1)a_k - \sum\limits_{k=1}^{n-1} (k-1)a_k}{n - (n-1)} \\
  & = \lim\limits_{n \to \infty} \dfrac{(n-1)a_n}{1} = 0.
\end{align*}
这表明 $\sum\limits_{n=1}^{\infty} a_n = \lim\limits_{n \to \infty} s_n = \lim\limits_{n \to \infty} \sigma_n = A.$ \qed

\end{frame}

%------------------------------------------------

\begin{frame}{更强的Tauber型定理}

实际上, 我们可以减弱 Tauber 定理中关于 $a_n = \mathrm{o} \left( \dfrac{1}{n} \right)$ 的假设, 得到一个更强的结果:

\begin{thm}
  若级数 $\sum\limits_{n=1}^{\infty} a_n = A~(c, 1),$ 且当 $n\to\infty$ 时 $a_n = \mathrm{O} \left( \dfrac{1}{n} \right),$ 那么该级数在通常意义下可和, 且 $\sum\limits_{n=1}^{\infty} a_n = A.$
\end{thm}

\vspace*{1em}

注意: 这里关于 $a_n$ 趋于 $0$ 速度的假设, 我们从 ${\color{red} \lim\limits_{n\to\infty} n a_n = 0},$ 减弱为了: 存在正实数 $M,$ 使得对任意自然数 $n$ 有 ${\color{red} \rvert n a_n \rvert \leqslant M}.$

\end{frame}

%------------------------------------------------

\begin{frame}{增强的Tauber型定理的分析与证明}

除了已知条件:当 $n\to\infty$ 时 $a_n = \mathrm{O} \left( \dfrac{1}{n} \right),$ 即存在正实数 $M,$ 使得对任意自然数 $n$ 有 ${\color{red} \rvert n a_n \rvert \leqslant M}$ 之外, 我们所拥有的条件是 $\sum\limits_{n=1}^{\infty} a_n = A~(c, 1),$ 即 $\lim\limits_{n\to\infty} \sigma_n = A,$ 这等价于 $\forall \varepsilon > 0,$ 存在自然数 $N$ 使得 ${\color{red} \forall n > m > N, \lvert \sigma_n - \sigma_m \vert < \varepsilon}.$

\pause
\vspace*{0.4em}

之前证明 Tauber 定理的时候, 我们导出了 $s_n$ 与 $\sigma_n$ 的关系
$$s_n = \sigma_n + \sum\limits_{k = 1}^n \dfrac{(k-1)a_k}{n}.$$
由于我们只有 $\rvert n a_n \rvert \leqslant M,$ 所以Stolz公式不能再用了, 得另想办法.

\pause
\vspace*{0.4em}

\begin{attnblock}
很多判别法、计算公式，例如 L'Hopital 法则, Stolz 公式, 包括后面要学的一致收敛的函数项级数的逐项微分、积分的结论, 相关条件都是一种充分性条件, 远\textbf{\color{red}不是必要条件}.
\end{attnblock}

\end{frame}

%------------------------------------------------

\begin{frame}{增强的Tauber型定理的分析与证明}

想法: 将 $s_n - \sigma_n$ 与 $\sigma_n - \sigma_m$ 联系上, 并用它控制余项的大小.

\only<2>{
假设 $n > m,$ 那么有
\begin{align*}
  \sigma_n - \sigma_m & = \dfrac{1}{n} \sum\limits_{k=1}^n s_k - \dfrac{1}{m} \sum\limits_{k=1}^{{\color{red}m}} s_k \\
  & = \dfrac{1}{n} \sum\limits_{k=1}^n s_k - \dfrac{1}{m} \sum\limits_{k=1}^{{\color{red}n}} s_k + \dfrac{1}{m} \sum\limits_{{\color{red}k=m+1}}^{n} s_k \\
  & = \dfrac{m - n}{nm} \sum\limits_{k=1}^n s_k + \dfrac{1}{m} \sum\limits_{k=m+1}^{n} s_k.
\end{align*}

另一方面, 有
$$s_n - \sigma_n = s_n - \dfrac{1}{n} \sum\limits_{k = 1}^n s_k.$$
}

\only<3>{假设 $n > m,$ 那么有
\begin{align*}
  \sigma_n - \sigma_m & = \dfrac{1}{n} \sum\limits_{k=1}^n s_k - \dfrac{1}{m} \sum\limits_{k=1}^m s_k \\
  & = \dfrac{1}{n} \sum\limits_{k=1}^n s_k - \dfrac{1}{m} \sum\limits_{k=1}^{n} s_k + \dfrac{1}{m} \sum\limits_{k=m+1}^{n} s_k \\
  & = \dfrac{m - n}{nm} {\color{red}
  \sum\limits_{k=1}^n s_k} + \dfrac{1}{m} \sum\limits_{k=m+1}^{n} s_k.
\end{align*}

另一方面, 有
$$s_n - \sigma_n = s_n - \dfrac{1}{n} {\color{red} \sum\limits_{k = 1}^n s_k}.$$
}

\end{frame}

%------------------------------------------------

\begin{frame}{增强的Tauber型定理的分析与证明}

两式通分相减, 有
\begin{align*}
& (s_n - \sigma_n) - \dfrac{m}{n - m}(\sigma_n - \sigma_m) \\
& = s_n - {\color{red} \cancel{\dfrac{1}{n} \sum\limits_{k = 1}^n s_k}} + {\color{red} \cancel{\dfrac{1}{n} \sum\limits_{k=1}^n s_k}} - \dfrac{1}{n-m} \sum\limits_{k=m+1}^{n} s_k \\
& = s_n - \dfrac{1}{n-m} \sum\limits_{k=m+1}^{n} s_k = \dfrac{1}{n-m} \sum\limits_{k=m+1}^{n} (s_n - s_k).
\end{align*}

\pause
\vspace*{1em}

整理一下, 有
$$s_n - \sigma_n = \underbrace{\color{cyan}\dfrac{m}{n - m}(\sigma_n - \sigma_m)}_{\text{\ding{172}}} + \underbrace{\color{magenta} \dfrac{1}{n-m} \sum\limits_{k=m+1}^{n} (s_n - s_k)}_{\text{\ding{173}}}.$$

\end{frame}

%------------------------------------------------

\begin{frame}{增强的Tauber型定理的分析与证明}

对于 $\text{\ding{173}} = \dfrac{1}{n-m} \sum\limits_{k=m+1}^{n} (s_n - s_k)$ 和式中的每一项, 我们有估计
{\smaller
\begin{multline*}
  \lvert s_n - s_k \rvert = \lvert a_{k+1} + \cdots + a_n \rvert \\
  \leqslant \dfrac{M}{k+1} + \cdots \dfrac{M}{n} \leqslant \dfrac{(n-k)M}{k+1} \leqslant \dfrac{(n-m-1)M}{m+1}
\end{multline*}
}
于是
{\smaller
$$
\left\lvert \dfrac{1}{n-m} \sum\limits_{k=m+1}^{n} (s_n - s_k) \right\rvert \leqslant \dfrac{1}{n-m} (n-m) \dfrac{n-m-1}{m+1} M = \dfrac{n-m-1}{m+1} M
$$
}

\pause

我们希望有 $\dfrac{n-m-1}{m+1} < \varepsilon$ 可推出 $m > \dfrac{n}{1 + \varepsilon} - 1.$ 取
$$\color{red} m = \left[ \dfrac{n}{1 + \varepsilon} \right] < n,$$
即有 $\lvert \text{\ding{173}} \rvert < M \varepsilon.$

\end{frame}

%------------------------------------------------

\begin{frame}{增强的Tauber型定理的分析与证明}

对于 $\text{\ding{172}} = \dfrac{m}{n - m}(\sigma_n - \sigma_m),$ 从 $m = \left[ \dfrac{n}{1 + \varepsilon} \right] \leqslant \dfrac{n}{1 + \varepsilon},$ 有 $m \varepsilon + m \leqslant n,$ 从而有 $\dfrac{m}{n - m} \leqslant \dfrac{1}{\varepsilon}.$ 于是
$$\lvert \text{\ding{172}} \rvert = \left\lvert \dfrac{m}{n - m}(\sigma_n - \sigma_m) \right\rvert \leqslant \dfrac{1}{\varepsilon} \left\lvert \sigma_n - \sigma_m \right\rvert.$$
若取 $N$ 充分大, 使得 $\forall n > m > N$ 都有 $\left\lvert \sigma_n - \sigma_m \right\rvert < {\color{red} \varepsilon^2},$ 那么
$$\left\lvert s_n - \sigma_n \right\rvert \leqslant \lvert \text{\ding{172}} \rvert + \lvert \text{\ding{173}} \rvert < \varepsilon + M \varepsilon = (M+1)\varepsilon.$$

\pause
\vspace{1em}

将上述思考、分析过程整理、调整顺序之后，即是完整的证明。

\end{frame}

%------------------------------------------------

\begin{frame}{Tauber型定理}

以上我们证明了两个所谓的Tauber型的定理, 即在某些附加的“正则性”条件下，根据考察的量的某些均值的性质，对这个量本身的性质作一些判断。

\pause
\vspace{1em}

这里, 我们附加了
$$a_n = s_n - s_{n-1} = \mathrm{o}\left(\dfrac{1}{n}\right) \text{ 或 } \mathrm{O}\left(\dfrac{1}{n}\right)$$
的“正则性”条件, 根据 $s_n$ 的均值 $\sigma_n$ 的收敛性, 得到了 $s_n$ 本身的收敛性.

\end{frame}

%------------------------------------------------

\section{总结}

%------------------------------------------------

\begin{frame}{总结}

\begin{itemize}
\item[{\larger[2] \color{red} \ding{43}}] Cesàro $(c, r)$ 意义下可和的级数 $\sum\limits_{n=1}^\infty a_n:$
\begin{align*}
  & (c, 0): \quad s_n = \dfrac{1}{n} \sum\limits_{k=1}^n a_n \text{ 收敛} \\
  & (c, 1): \quad \sigma_n = \dfrac{1}{n} \sum\limits_{k=1}^n s_n \text{ 收敛} \\
  & \phantom{++++++} \vdots
\end{align*}
\vspace{-0.8em}
\pause
\item[{\larger[2] \color{red} \ding{43}}] 有集合真包含关系: $\{ (c, r) \text{可和级数} \} \subsetneqq \{ (c, r + 1) \text{可和级数} \}$
\pause
\vspace{0.8em}
\item[{\larger[2] \color{red} \ding{43}}] Tauber型定理:
$$
\left. \begin{array}{r}
\sum\limits_{n=1}^\infty a_n \text{ 在 $(c,1)$ 意义下可和} \\
a_n = \mathrm{o} \left( \dfrac{1}{n} \right) (\text{或 } \mathrm{O} \left( \dfrac{1}{n} \right))
\end{array} \right\} \Longrightarrow \sum\limits_{n=1}^\infty a_n \text{ 在通常意义下可和}
$$
\end{itemize}

\end{frame}

%------------------------------------------------

%------------------------------------------------
% % closing page
% \begin{frame}
% \HUGE{\centerline{\bfseries {\color{gray} The} {\color{zkblue} End}}}

% \pause

% \vspace{1.2em}

% \Large{\centerline{\bfseries {\color{gray}Thank} \quad {\color{zkblue} You}}}

% \end{frame}

%----------------------------------------------------------------------------------------

\end{document}
